%%%%%%%%%%%%%%%%%%%%%%%%%%%%%%%%%%%%%%%%%
% Lachaise Assignment
% LaTeX Template
% Version 1.0 (26/6/2018)
%
% This template originates from:
% http://www.LaTeXTemplates.com
%
% Authors:
% Marion Lachaise & François Févotte
% Vel (vel@LaTeXTemplates.com)
%
% License:
% CC BY-NC-SA 3.0 (http://creativecommons.org/licenses/by-nc-sa/3.0/)
% 
%%%%%%%%%%%%%%%%%%%%%%%%%%%%%%%%%%%%%%%%%

%----------------------------------------------------------------------------------------
%	PACKAGES AND OTHER DOCUMENT CONFIGURATIONS
%----------------------------------------------------------------------------------------

\documentclass{article}
\usepackage{amsthm}
\usepackage{listings}
\usepackage{xcolor}
\usepackage{graphicx}
\usepackage{cite}
\usepackage{longtable}
\usepackage{array}
\usepackage{booktabs}

\lstset{
	language=Python,
	basicstyle=\ttfamily\small,
	keywordstyle=\color{blue},
	stringstyle=\color{red},
	commentstyle=\color{purple},
	showstringspaces=false,
	numbers=left,
	numberstyle=\tiny\color{gray},
	breaklines=true,
	frame=single,
	captionpos=b
}
\newtheorem{definition}{Definition}[section]

%%%%%%%%%%%%%%%%%%%%%%%%%%%%%%%%%%%%%%%%%
% Lachaise Assignment
% Structure Specification File
% Version 1.0 (26/6/2018)
%
% This template originates from:
% http://www.LaTeXTemplates.com
%
% Authors:
% Marion Lachaise & François Févotte
% Vel (vel@LaTeXTemplates.com)
%
% License:
% CC BY-NC-SA 3.0 (http://creativecommons.org/licenses/by-nc-sa/3.0/)
% 
%%%%%%%%%%%%%%%%%%%%%%%%%%%%%%%%%%%%%%%%%

%----------------------------------------------------------------------------------------
%	PACKAGES AND OTHER DOCUMENT CONFIGURATIONS
%----------------------------------------------------------------------------------------

\usepackage{amsmath,amsfonts,stmaryrd,amssymb} % Math packages

\usepackage{enumerate} % Custom item numbers for enumerations

\usepackage[ruled]{algorithm2e} % Algorithms

\usepackage[framemethod=tikz]{mdframed} % Allows defining custom boxed/framed environments

\usepackage{listings} % File listings, with syntax highlighting
\lstset{
	basicstyle=\ttfamily, % Typeset listings in monospace font
}

%----------------------------------------------------------------------------------------
%	DOCUMENT MARGINS
%----------------------------------------------------------------------------------------

\usepackage{geometry} % Required for adjusting page dimensions and margins

\geometry{
	paper=a4paper, % Paper size, change to letterpaper for US letter size
	top=2.5cm, % Top margin
	bottom=3cm, % Bottom margin
	left=2.5cm, % Left margin
	right=2.5cm, % Right margin
	headheight=14pt, % Header height
	footskip=1.5cm, % Space from the bottom margin to the baseline of the footer
	headsep=1.2cm, % Space from the top margin to the baseline of the header
	%showframe, % Uncomment to show how the type block is set on the page
}

%----------------------------------------------------------------------------------------
%	FONTS
%----------------------------------------------------------------------------------------

\usepackage[utf8]{inputenc} % Required for inputting international characters
\usepackage[T1]{fontenc} % Output font encoding for international characters

\usepackage{XCharter} % Use the XCharter fonts

%----------------------------------------------------------------------------------------
%	COMMAND LINE ENVIRONMENT
%----------------------------------------------------------------------------------------

% Usage:
% \begin{commandline}
%	\begin{verbatim}
%		$ ls
%		
%		Applications	Desktop	...
%	\end{verbatim}
% \end{commandline}

\mdfdefinestyle{commandline}{
	leftmargin=10pt,
	rightmargin=10pt,
	innerleftmargin=15pt,
	middlelinecolor=black!50!white,
	middlelinewidth=2pt,
	frametitlerule=false,
	backgroundcolor=black!5!white,
	frametitle={Command Line},
	frametitlefont={\normalfont\sffamily\color{white}\hspace{-1em}},
	frametitlebackgroundcolor=black!50!white,
	nobreak,
}

% Define a custom environment for command-line snapshots
\newenvironment{commandline}{
	\medskip
	\begin{mdframed}[style=commandline]
}{
	\end{mdframed}
	\medskip
}

%----------------------------------------------------------------------------------------
%	FILE CONTENTS ENVIRONMENT
%----------------------------------------------------------------------------------------

% Usage:
% \begin{file}[optional filename, defaults to "File"]
%	File contents, for example, with a listings environment
% \end{file}

\mdfdefinestyle{file}{
	innertopmargin=1.6\baselineskip,
	innerbottommargin=0.8\baselineskip,
	topline=false, bottomline=false,
	leftline=false, rightline=false,
	leftmargin=2cm,
	rightmargin=2cm,
	singleextra={%
		\draw[fill=black!10!white](P)++(0,-1.2em)rectangle(P-|O);
		\node[anchor=north west]
		at(P-|O){\ttfamily\mdfilename};
		%
		\def\l{3em}
		\draw(O-|P)++(-\l,0)--++(\l,\l)--(P)--(P-|O)--(O)--cycle;
		\draw(O-|P)++(-\l,0)--++(0,\l)--++(\l,0);
	},
	nobreak,
}

% Define a custom environment for file contents
\newenvironment{file}[1][File]{ % Set the default filename to "File"
	\medskip
	\newcommand{\mdfilename}{#1}
	\begin{mdframed}[style=file]
}{
	\end{mdframed}
	\medskip
}

%----------------------------------------------------------------------------------------
%	NUMBERED QUESTIONS ENVIRONMENT
%----------------------------------------------------------------------------------------

% Usage:
% \begin{question}[optional title]
%	Question contents
% \end{question}

\mdfdefinestyle{question}{
	innertopmargin=1.2\baselineskip,
	innerbottommargin=0.8\baselineskip,
	roundcorner=5pt,
	nobreak,
	singleextra={%
		\draw(P-|O)node[xshift=1em,anchor=west,fill=white,draw,rounded corners=5pt]{%
		Question \theQuestion\questionTitle};
	},
}

\newcounter{Question} % Stores the current question number that gets iterated with each new question

% Define a custom environment for numbered questions
\newenvironment{question}[1][\unskip]{
	\bigskip
	\stepcounter{Question}
	\newcommand{\questionTitle}{~#1}
	\begin{mdframed}[style=question]
}{
	\end{mdframed}
	\medskip
}

%----------------------------------------------------------------------------------------
%	WARNING TEXT ENVIRONMENT
%----------------------------------------------------------------------------------------

% Usage:
% \begin{warn}[optional title, defaults to "Warning:"]
%	Contents
% \end{warn}

\mdfdefinestyle{warning}{
	topline=false, bottomline=false,
	leftline=false, rightline=false,
	nobreak,
	singleextra={%
		\draw(P-|O)++(-0.5em,0)node(tmp1){};
		\draw(P-|O)++(0.5em,0)node(tmp2){};
		\fill[black,rotate around={45:(P-|O)}](tmp1)rectangle(tmp2);
		\node at(P-|O){\color{white}\scriptsize\bf !};
		\draw[very thick](P-|O)++(0,-1em)--(O);%--(O-|P);
	}
}

% Define a custom environment for warning text
\newenvironment{warn}[1][Warning:]{ % Set the default warning to "Warning:"
	\medskip
	\begin{mdframed}[style=warning]
		\noindent{\textbf{#1}}
}{
	\end{mdframed}
}

%----------------------------------------------------------------------------------------
%	INFORMATION ENVIRONMENT
%----------------------------------------------------------------------------------------

% Usage:
% \begin{info}[optional title, defaults to "Info:"]
% 	contents
% 	\end{info}

\mdfdefinestyle{info}{%
	topline=false, bottomline=false,
	leftline=false, rightline=false,
	nobreak,
	singleextra={%
		\fill[black](P-|O)circle[radius=0.4em];
		\node at(P-|O){\color{white}\scriptsize\bf i};
		\draw[very thick](P-|O)++(0,-0.8em)--(O);%--(O-|P);
	}
}

% Define a custom environment for information
\newenvironment{info}[1][Info:]{ % Set the default title to "Info:"
	\medskip
	\begin{mdframed}[style=info]
		\noindent{\textbf{#1}}
}{
	\end{mdframed}
}
 % Include the file specifying the document structure and custom commands

%----------------------------------------------------------------------------------------
%	ASSIGNMENT INFORMATION
%----------------------------------------------------------------------------------------

\title{Tarefa IV PCS5028 - Metodologia de Pesquisa Científica em Engenharia de Computação} % Title of the assignment

\author{Giovani Tavares\\ \texttt{giovanitavares@usp.br}} % Author name and email address

\date{Universidade de São Paulo --- July 12, 2026} % University, school and/or department name(s) and a date

%----------------------------------------------------------------------------------------

\begin{document}

\maketitle % Print the title

%----------------------------------------------------------------------------------------
%	INTRODUCTION
%----------------------------------------------------------------------------------------





\section{Resumo}
% =====================================================================

Sistemas de \textit{Retrieval-Augmented Generation} (RAG) combinam grandes modelos de linguagem (LLMs) com mecanismos de recuperação de informação. Arquiteturas baseadas exclusivamente em similaridade vetorial apresentam limitações conhecidas na literatura como a recuperação de contextos semanticamente próximos, porém nem sempre factualmente pertinentes, comprometendo métricas como \textit{faithfulness} e \textit{context recall}.
Nesse sentido, trabalhos recentes relatam ganhos consistentes em arquiteturas que integram uma camada de representação simbólica do conhecimento (ontologias e grafos) no processo de recuperação do RAG, propondo sistemas que se incluem no que se denomina IA Neurossimbólica. O OG RAG \cite{sharma2025ograg} demonstrou tais ganhos nos domínios de agricultura e jornalismo, enquanto o \cite{feng2025ontologyrag} o fez no  contexto da biomedicina. Há ainda aplicações nos contextos jurídico-regulatórias e \textit{compliance} financeiro. Este projeto investiga se o mesmo efeito se replica no domínio de modelos metabólicos em escala genômica (\textit{genome-scale metabolic models}, GEMs), propondo um estudo comparativo controlado entre três arquiteturas: RAG vetorial puro, GraphRAG e OG-RAG. Elas serão instanciadas sobre a ontologia de domínio OntoGEM \cite{souza2025ontogem} e sobre dados do repositório BiGG Models, com foco no modelo \textit{E. coli core}. A hipótese central é que a inclusão da camada simbólica melhora as métricas do \textit{framework} RAGAS \cite{es-etal-2024-ragas} neste domínio. Para testá-la, será construído, com apoio de especialistas, um conjunto anotado de perguntas de competência sobre o modelo metabólico de \textit{E. coli}, pré-classificadas por tipo de explicação exigida, com respostas de referência derivadas de forma reprodutível a partir da própria base BiGG. As três arquiteturas serão avaliadas sob condições controladas (mesmo LLM gerador, mesmo corpus-fonte, variando apenas o mecanismo de recuperação), isolando o efeito da estrutura simbólica sobre as métricas de \textit{faithfulness}, \textit{context precision}, \textit{context recall} e \textit{answer relevancy}. Espera-se produzir evidência empírica, ainda que restrita a este domínio e modelo, sobre o papel de camadas de conhecimento auditáveis e editáveis por especialistas na melhoria de sistemas de perguntas e respostas baseados em RAG.

\newpage

% =====================================================================
\section{Introdução e Justificativa}
% =====================================================================

\subsection{Contextualização}

Sistemas de \textit{Retrieval-Augmented Generation} \cite{lewis2020retrieval} tornaram-se a abordagem dominante para reduzir alucinações e ancorar respostas de LLMs em fontes verificáveis, recuperando trechos de um corpus externo por similaridade vetorial e condicionando a geração da resposta a esse contexto. Essa abordagem, no entanto, herda uma limitação estrutural: a similaridade vetorial captura proximidade semântica, não correção factual ou relevância segundo a estrutura do domínio, de modo que o contexto recuperado pode ser semanticamente próximo da pergunta e, ainda assim, irrelevante ou mesmo contraditório em relação aos fatos que os especialistas de domínio estabelecem como corretos. 

Buscando mitigar essa limitação, a integração de representações simbólicas do conhecimento validadas por especialistas a arquiteturas de RAG tem-se demonstrado promissora em diferentes domínios \cite{sharma2025ograg} \cite{feng2025ontologyrag}. Os sistemas nascentes de tal integração fazem parte do que se denomina Inteligência Artifical Neurossimbólica. Além das melhorias em métricas como \textit{faithfulness}, \textit{context precision}, \textit{context recall} e \textit{answer relevancy}, sistemas neurossimbólicos como o OG RAG \cite{sharma2025ograg} são dotados de explicabilidade e rastreabilidade estrutural e não somente estatística como o RAG com somente busca vetorial. A integração de ontologias ao processo de recuperação de informação do RAG proposta em \cite{sharma2025ograg} permite que cada fato recuperado carregue a relação semântica que o liga à pergunta, enquanto que a busca vetorial pura apenas relaciona o fato à pergunta com métricas de similaridade vetorial de pouca interpretabilidade por seres humanos.

\subsection{Bibliografia}

O \textit{framework} RAGAS \cite{es-etal-2024-ragas} formalizou a avaliação de pipelines de RAG por meio de métricas derivadas automaticamente da tripla pergunta-resposta-contexto, entre as quais se destacam \textit{faithfulness} (grau em que a resposta é sustentada pelo contexto recuperado), \textit{context precision} e \textit{context recall} (adequação e completude do contexto recuperado em relação à resposta de referência) e \textit{answer relevancy}. Por não depender de anotação humana exaustiva a cada execução, o RAGAS tornou-se referência para comparação de arquiteturas de RAG e é adotado, direta ou indiretamente, pelos trabalhos discutidos a seguir.

Sharma et al. propõem o OG-RAG (\textit{Ontology-Grounded RAG}) \cite{sharma2025ograg}, no qual uma ontologia de domínio é usada como esquema de extração e recuperação (o que os autores chamam de Modo A), estruturando o corpus em fatos compatíveis com a ontologia antes da indexação; os autores relatam ganhos em métricas inspiradas no RAGAS nos domínios de agricultura e jornalismo. Feng et al. \cite{feng2025ontologyrag} aplicam uma estratégia análoga ao mapeamento de códigos biomédicos, ancorando a recuperação em hierarquias terminológicas controladas (por exemplo, SNOMED CT e ICD), com ganhos reportados principalmente em \textit{faithfulness}. Kalra et al. \cite{kalra2024hyparag} propõem uma arquitetura híbrida para aplicações jurídicas e regulatórias, combinando grafos de conhecimento hierárquicos com recuperação vetorial, relatando melhorias em \textit{context recall} frente a baselines puramente vetoriais. Nie et al. \cite{nie2025aml} aplicam integração simbólica a \textit{compliance} financeiro (\textit{anti-money laundering}), reportando reduções de contexto espúrio quando a recuperação é restringida por regras e relações de domínio validadas.

O GraphRAG \cite{edge2024graphrag} representa uma outra  família de arquiteturas neurossimbólicas, na qual a estrutura de conhecimento não é uma ontologia curada por especialistas, mas um grafo de entidades e relações extraído automaticamente do corpus por um LLM, posteriormente utilizado para sumarização hierárquica e recuperação. Essa distinção é relevante para o presente projeto: GraphRAG e OG-RAG compartilham o uso de uma camada estruturada intermediária entre o corpus e a geração, mas diferem quanto à origem e à validação dessa estrutura --- extraída automaticamente no primeiro caso, curada e validada por especialistas de domínio no segundo. Comparar as duas arquiteturas contra um baseline vetorial permite não apenas testar se \emph{alguma} estrutura simbólica melhora as métricas do RAGAS, mas também explorar, de forma exploratória, se a origem e o grau de curadoria dessa estrutura importam.

A camada simbólica curada utilizada neste é a ontologia OntoGEM, desenvolvida por Souza \cite{souza2025ontogem} em seu projeto de doutorado do mesmo grupo de pesquisa. O trabalho de Souza propõe uma ontologia de domínio específica para Modelos Metabólicos de Escala Genômica (GEMs), formalizada a partir de perguntas de competência formadas junto a especialistas de biologia de sistemas e estruturada segundo etapas clássicas de engenharia de ontologias: especificação, conceitualização, formalização e implementação. O objetivo original dessa ontologia é acelerar e melhorar a qualidade do processo de reconstrução de GEMs, viabilizando a integração e reconciliação de dados provenientes de múltiplas bases biológicas (como BiGG, KEGG e ChEBI). Embora a motivação original da OntoGEM seja a reconstrução de modelos metabólicos, e não a recuperação de informação em sistemas de RAG, sua formalização das entidades e relações de um GEM (reações, metabólitos, genes, complexos enzimáticos, compartimentos celulares) a torna diretamente reaproveitável para a arquitetura OG-RAG neste projeto, dispensando a necessidade de desenvolver uma nova ontologia de domínio do zero.

Em conjunto, esses trabalhos sugerem um padrão consistente de ganhos em \textit{faithfulness} e \textit{context recall} associados à presença de estrutura simbólica no pipeline de recuperação, em diferentes domínios. Entretanto, nenhum desses domínios possui a característica central do domínio de modelos metabólicos em escala genômica. Não há, até onde este projeto pôde levantar, evidência publicada sobre o efeito de camadas simbólicas em RAG aplicado a esse domínio.


\subsection{O Domínio de Aplicação: Modelos Metabólicos em Escala Genômica}


O processo de construção de um modelo computacional do metabolismo de um organismo é conhecido como reconstrução, enquanto o modelo resultante é frequentemente denominado modelo metabólico em escala genômica (GEM) \cite{souza2025ontogem}. O repositório BiGG Models \cite{king2016bigg} disponibiliza GEMs curados e padronizados, entre os quais o modelo \textit{E. coli core}, um subconjunto didático e amplamente utilizado do metabolismo central de \textit{Escherichia coli}, adequado à construção de um conjunto controlado de perguntas de domínio por combinar cobertura representativa das entidades do modelo (reações, metabólitos, genes, compartimentos) com um tamanho factível para anotação por especialistas durante o período de mestrado.

A ontologia OntoGEM, desenvolvida em \cite{souza2025ontogem}, formaliza as entidades e relações de um GEM, sendo elas as reações, metabólitos, genes, complexos enzimáticos, compartimentos celulares e as relações estequiométricas e regulatórias entre eles. Dessa forma, pretende-se utilizar essa ontologia como o esquema de extração e recuperação da arquitetura OG-RAG.

\subsection{Lacuna e Motivação}

Não há evidência de que os ganhos reportados para RAG neurossimbólico em diferentes domínios também se aplicam ao domínio de perguntas e respostas relacionados a GEMs. Além disso, os trabalhos revisados comparam arquiteturas neurossimbólicas a baselines puramente vetoriais, mas raramente isolam o efeito específico da camada simbólica de outras variáveis do pipeline (modelo gerador, corpus, estratégia de \textit{chunking}), o que limita a atribuição causal dos ganhos observados. Este projeto propõe um desenho experimental controlado que mantém fixos o LLM gerador e o corpus-fonte, variando apenas o mecanismo de recuperação (vetorial puro, GraphRAG, OG-RAG), de modo a isolar o efeito da representação simbólica do conhecimento sobre as métricas do RAGAS.

\newpage

% =====================================================================
\section{Objetivos}
% =====================================================================

\subsection{Objetivo Geral}

Demonstrar o efeito da inclusão da representação simbólica do conhecimento em diferentes arquiteturas de RAG sobre as métricas do RAGAS \cite{es-etal-2024-ragas}, no domínio de modelos metabólicos em escala genômica.

\subsection{Objetivos Específicos}

\begin{enumerate}
	\item Construir, com apoio de especialistas de domínio, um conjunto anotado de perguntas de competência sobre o modelo metabólico \textit{E. coli core}, pré-classificado por tipo de explicação exigida, com respostas de referência derivadas de forma reprodutível a partir do repositório BiGG Models.
	\item Formalizar a ontologia OntoGEM como esquema de extração e recuperação (Modo A) para uso nas arquiteturas OG-RAG.
	\item Implementar três pipelines de RAG comparáveis --- vetorial puro (\textit{baseline}), GraphRAG e OG-RAG --- operando sobre o mesmo corpus-fonte e o mesmo modelo gerador, variando apenas o mecanismo de recuperação.
	\item Avaliar as três arquiteturas segundo as métricas do RAGAS (\textit{faithfulness}, \textit{context precision}, \textit{context recall} e \textit{answer relevancy}) sobre o conjunto de perguntas construído.
	\item Analisar estatisticamente se as diferenças observadas entre arquiteturas com e sem camada simbólica são significativas, e se o padrão de ganhos (concentração em \textit{faithfulness} e \textit{context recall}) reportado na literatura de outros domínios se replica no domínio de GEMs.
	\item Discutir, a partir dos resultados, os limites de generalização das evidências obtidas para modelos e domínios distintos do estudado.
\end{enumerate}

\newpage

% =====================================================================
\section{Plano de Trabalho e Cronograma de Execução}
% =====================================================================

O projeto está organizado em seis fases, detalhadas a seguir, executadas ao longo de 24 meses.

\begin{enumerate}
	\item \textbf{Revisão bibliográfica.} Consolidação da revisão sistemática sobre RAG neurossimbólico, arquiteturas de recuperação estruturada (OG-RAG, GraphRAG) e literatura de avaliação (RAGAS).
	\item \textbf{Refinamento da OntoGEM para os frameworks de RAG e exame de qualificação.} Adaptação técnica da ontologia OntoGEM \cite{souza2025ontogem} --- cuja formalização conceitual é resultado do trabalho de Souza \cite{souza2025ontogem} --- para os formatos e esquemas exigidos pelos \textit{frameworks} de RAG utilizados (por exemplo, conversão da OWL para o esquema de extração/recuperação do OG-RAG Modo A), sem alterar sua estrutura conceitual original. Em paralelo, realização do exame de qualificação.
	\item \textbf{Construção do dataset anotado de perguntas de domínio e implementação das arquiteturas de RAG.} Elaboração, em conjunto com especialistas em biologia de sistemas/metabolismo, de um conjunto de perguntas de competência sobre o modelo \textit{E. coli core}, com classificação por tipo de explicação e construção de respostas de referência (detalhado na Seção 5.2). Em paralelo, desenvolvimento dos três pipelines de RAG (vetorial puro, GraphRAG, OG-RAG Modo A) sobre corpus comum derivado do BiGG Models e da OntoGEM adaptada.
	\item \textbf{Execução dos experimentos e coleta de métricas.} Execução das três arquiteturas sobre o conjunto de perguntas, coleta das métricas RAGAS e dos artefatos intermediários (contextos recuperados, respostas geradas).
	\item \textbf{Análise dos resultados, redação e defesa.} Análise estatística comparativa entre arquiteturas, análise por tipo de pergunta, redação da dissertação e defesa final.
\end{enumerate}

\begin{table}[h]
	\centering
	\resizebox{\textwidth}{!}{%
		\begin{tabular}{@{}lcccccccc@{}}
			\toprule
			\textbf{Atividade} & \textbf{T1} & \textbf{T2} & \textbf{T3} & \textbf{T4} & \textbf{T5} & \textbf{T6} & \textbf{T7} & \textbf{T8} \\
			\midrule
			Revisão bibliográfica & X & X & & & & & & \\
			Refinamento da OntoGEM para os frameworks de RAG & & & X & X & & & & \\
			Exame de qualificação & & & X & X & & & & \\
			Construção do dataset anotado (especialistas) & & & & X & X & & & \\
			Implementação das arquiteturas de RAG (vetorial, GraphRAG, OG-RAG) & & & & X & X & & & \\
			Execução dos experimentos e coleta de resultados (RAGAS) & & & & & & X & & \\
			Análise estatística dos resultados & & & & & & & X & \\
			Redação da dissertação & & & & & & & X & X \\
			Defesa da dissertação & & & & & & & & X \\
			\bottomrule
		\end{tabular}%
	}
\end{table}

\noindent \textit{Nota: T1--T8 referem-se a trimestres consecutivos ao longo de 24 meses.}

\newpage

% =====================================================================
\section{Material e Métodos}
% =====================================================================

\subsection{Desenho Experimental}

O estudo adota um desenho comparativo de três fases, no qual a variável manipulada é exclusivamente o mecanismo de recuperação, mantendo-se fixos o LLM gerador, o corpus-fonte e o conjunto de perguntas de avaliação, de modo a isolar o efeito da representação simbólica do conhecimento sobre as métricas do RAGAS:

\begin{itemize}
	\item \textbf{Fase A --- RAG vetorial (baseline).} Indexação do corpus textual derivado do BiGG Models por \textit{embeddings} densos, com recuperação por similaridade de cosseno, sem qualquer estrutura simbólica intermediária. Representa o estado da arte sem integração de conhecimento estruturado.
	\item \textbf{Fase B --- GraphRAG \cite{edge2024graphrag}.} Extração automática, via LLM, de um grafo de entidades e relações a partir do mesmo corpus, com sumarização hierárquica de comunidades e recuperação orientada pelo grafo. A estrutura simbólica não é curada por especialistas, mas induzida a partir dos dados.
	\item \textbf{Fase C --- OG-RAG, Modo A \cite{sharma2025ograg}.} Uso da ontologia OntoGEM \cite{souza2025ontogem} como esquema de extração e recuperação: o corpus é estruturado em fatos compatíveis com as classes e relações da ontologia antes da indexação, e a recuperação é restringida pelo esquema ontológico. Única fase com estrutura simbólica curada e validada por especialistas de domínio (Modo A, conforme operacionalização de \cite{sharma2025ograg}; os Modos B --- vocabulário de \textit{prompt} --- e C --- validador \textit{a posteriori} --- não são objeto deste estudo).
\end{itemize}

As três fases utilizarão o mesmo modelo gerador (LLM), a mesma estratégia de \textit{prompt} de geração de resposta final (variando apenas o contexto recuperado) e o mesmo corpus-fonte, derivado do modelo \textit{E. coli core} do repositório BiGG Models.

\subsection{Construção do Dataset Anotado de Perguntas de Domínio}

Um dos pilares metodológicos do projeto é a construção, junto a especialistas de domínio (pesquisadores da área de biologia de sistemas/modelagem metabólica), de um conjunto anotado de perguntas de competência sobre o modelo \textit{E. coli core}. Essa etapa seguirá o seguinte protocolo:

\begin{enumerate}
	\item \textbf{Reaproveitamento das perguntas de competência.} As perguntas de competência utilizadas neste projeto são as já elicitadas e formuladas no âmbito do trabalho de Souza \cite{souza2025ontogem}, que cobrem consultas sobre metabólitos, reações e topologia da rede (por exemplo, estequiometria de reações, direcionalidade, genes associados via regras GPR, compartimentalização e vias metabólicas). Não há, portanto, uma nova etapa de elicitação de perguntas junto a especialistas neste projeto.
	\item \textbf{Seleção das reações para instanciação das perguntas.} Diversas perguntas de competência são parametrizadas em termos de uma reação genérica $R$ (ou de outras entidades genéricas do modelo), exigindo a seleção de instâncias concretas do modelo \textit{E. coli core} para gerar as perguntas efetivamente utilizadas na avaliação. Essa seleção considerará (i) a dificuldade esperada da pergunta em função da reação escolhida (por exemplo, reações com múltiplas isoformas enzimáticas ou regras GPR complexas tendem a gerar perguntas mais difíceis do que reações com associação gene-reação simples) e (ii) a quantidade de instanciações necessária para garantir cobertura representativa das entidades e relações do modelo em cada categoria de pergunta.
	\item \textbf{Reaproveitamento da classificação por tipo de explicação.} As perguntas de competência de Souza \cite{souza2025ontogem} já são acompanhadas de uma classificação por categoria de explicação, definida independentemente da execução deste projeto, o que evita classificações \textit{post hoc} e o viés de confirmação associado. As categorias reaproveitadas são: \textit{reaction provenance}, \textit{metabolite provenance}, \textit{model provenance}, \textit{biological evidence}, \textit{model assumptions}, \textit{model validation}, \textit{network structure}, \textit{semantic classification}, \textit{conflict management} e \textit{accountability}. Essa classificação será utilizada diretamente como base para a análise estratificada dos resultados por categoria.
	\item \textbf{Construção de respostas de referência (\textit{gold answers}).} As respostas de referência serão construídas em duas etapas: (i) busca e extração dos dados correspondentes diretamente na API do BiGG Models para cada instância selecionada, garantindo rastreabilidade entre cada resposta de referência e o dado-fonte correspondente; (ii) sessões de validação com especialistas de domínio para conferência e, quando necessário, correção das respostas obtidas.
	\item \textbf{Pré-registro de hipóteses por categoria.} Dentre as categorias de pergunta reaproveitadas de Souza \cite{souza2025ontogem}, aquelas nas quais se espera maior efeito da camada simbólica (por exemplo, \textit{network structure}, \textit{reaction provenance} e \textit{metabolite provenance}, que dependem mais fortemente de restrições estruturais do modelo) serão declaradas antes da execução dos experimentos, permitindo uma análise confirmatória dessas categorias específicas, além da análise exploratória das outras.
\end{enumerate}

O conjunto resultante --- da ordem de dezenas de perguntas de competência, dimensionado para cobertura representativa das entidades e relações do modelo \textit{E. coli core} --- constitui o instrumento de avaliação comum às três fases experimentais.

\subsection{Corpus e Recursos de Domínio}

\begin{itemize}
	\item \textbf{Fonte de dados:} modelo \textit{E. coli core}, obtido do repositório BiGG Models \cite{king2016bigg}, contendo reações, metabólitos, genes e regras de associação gene-proteína-reação (GPR).
	\item \textbf{Ontologia de domínio:} OntoGEM, desenvolvida em \cite{souza2025ontogem}, formalizando classes e relações estequiométricas e regulatórias do modelo, utilizada como esquema de extração/recuperação na Fase C.
	\item \textbf{Corpus textual comum:} representação em linguagem natural (formato \textit{markdown}) do modelo \textit{E. coli core}, gerada de forma determinística a partir dos dados estruturados do BiGG Models, garantindo que os três braços operem sobre o mesmo conteúdo informacional.
	
\end{itemize}

\subsection{Implementação}

As três arquiteturas serão implementadas com bibliotecas de código aberto de referência para cada abordagem (RAG vetorial, GraphRAG e OG-RAG), com o mesmo modelo de linguagem utilizado como gerador em todas as fases, de modo a evitar que diferenças de capacidade do LLM confundam o efeito da variável de interesse. Parâmetros comuns entre arquiteturas (tamanho de \textit{chunk}, número de documentos recuperados \textit{top-k}, temperatura de geração) serão mantidos constantes sempre que a arquitetura permitir.

\newpage

% =====================================================================
\section{Forma de Análise dos Resultados}
% =====================================================================

\subsection{Métricas}

Para cada par (pergunta, arquitetura), serão calculadas as métricas do RAGAS \cite{es-etal-2024-ragas}: \textit{faithfulness}, \textit{context precision}, \textit{context recall} e \textit{answer relevancy}. As distribuições dessas métricas serão comparadas entre as três fases experimentais.

\subsection{Testes Estatísticos}

Como cada pergunta do conjunto de avaliação criado será respondida pelas três arquiteturas, a comparação entre elas será feita a partir da diferença de cada métrica, pergunta a pergunta, entre uma arquitetura com camada simbólica (GraphRAG ou OG-RAG) e o baseline vetorial. Para testar se essas diferenças são estatisticamente significativas, será utilizado o teste de Wilcoxon signed-rank, uma alternativa não paramétrica ao teste t pareado. A escolha por um teste não paramétrico se justifica porque as métricas do RAGAS são limitadas ao intervalo $[0,1]$, o que tende a produzir distribuições assimétricas ou concentradas nas extremidades, violando a suposição de normalidade exigida pelos testes paramétricos. Como o p-valor depende fortemente do tamanho da amostra e não indica, por si só, a magnitude da diferença observada, os resultados dos testes serão complementados por uma medida de tamanho de efeito (a correlação \textit{rank-biserial} pareada), permitindo distinguir diferenças estatisticamente significativas mas pouco relevantes na prática de diferenças de fato substanciais entre arquiteturas. Essa distinção é fundamental neste trabalho dado o tamanho moderado do conjunto de perguntas de avaliação que será criado.

\subsection{Análise Estratificada por Tipo de Pergunta}

Consistente com o pré-registro das categorias de pergunta (Seção 5.2-e), os resultados serão analisados também de forma estratificada por tipo de explicação exigida, permitindo testar especificamente a hipótese de que o efeito da camada simbólica é mais pronunciado em perguntas relacionais/multi-\textit{hop}, nas quais a estrutura da ontologia atua como restrição sobre o espaço de recuperação.

\subsection{Análise Qualitativa}

Uma amostra de casos em que as arquiteturas divergem substancialmente em \textit{faithfulness} ou \textit{context recall} será inspecionada qualitativamente, com o objetivo de caracterizar os mecanismos pelos quais a estrutura simbólica altera o contexto recuperado.

\subsection{Critérios de Confirmação/Refutação da Hipótese}

A hipótese de que a inclusão da camada de representação simbólica do conhecimento melhora as métricas do RAGAS no domínio de GEMs será considerada válida se as fases com estrutura simbólica (GraphRAG e/ou OG-RAG) apresentarem ganhos estatisticamente significativos frente ao baseline vetorial, concentrados nas métricas de \textit{faithfulness} e \textit{context recall}, replicando o padrão relatado na literatura de outros domínios \cite{sharma2025ograg,feng2025ontologyrag,kalra2024hyparag,nie2025aml}. A comparação entre as Fases B e C permitirá discutir se o grau de curadoria da estrutura simbólica (induzida no GraphRAG versus validada por especialistas em OG RAG) influencia a magnitude dos ganhos observados.

\subsection{Resultados Esperados}

Espera-se demonstrar que a representação simbólica do conhecimento é responsável pela melhoria de métricas relacionadas a sistemas de perguntas e respostas baseados em RAG, fornecendo evidência sobre a importância de integrar camadas de conhecimento validadas, compreensíveis, auditáveis e editáveis por especialistas de domínio em arquiteturas até então baseadas exclusivamente em aprendizado de máquina --- das quais o RAG por busca vetorial pura é o principal exemplo.

\subsection{Limitações}

As evidências de melhoria, quando e se coletadas, estarão restritas às perguntas de domínio relacionadas ao modelo metabólico \textit{E. coli core}, selecionadas com apoio de especialistas de domínio. Melhorias observadas para este modelo e domínio específicos não devem ser generalizadas, sem validação adicional, a outros modelos metabólicos em escala genômica ou a outros domínios de aplicação, os quais estão fora do escopo desta pesquisa.


\bibliographystyle{plain}
\bibliography{references}

\end{document}
